\documentclass[12pt]{article}

% Idioma y codificación
\usepackage[spanish, es-tabla]{babel}       %es-tabla para que se titule "Tabla"
\usepackage[utf8]{inputenc}

% Márgenes
\usepackage[a4paper,top=3cm,bottom=2.5cm,left=3cm,right=3cm]{geometry}

% Comentarios de bloque
\usepackage{verbatim}

% Paquetes de links
\usepackage[hidelinks]{hyperref}    % Permite enlaces
\usepackage{url}                    % redirecciona a la web

% Más opciones para enumeraciones
\usepackage{enumitem}

% Personalizar la portada
\usepackage{titling}


% Paquetes de tablas
\usepackage{multirow}


%------------------------------------------------------------------------

%Paquetes de figuras
\usepackage{caption}
\usepackage{subcaption} % Figuras al lado de otras
\usepackage{float}      % Poner figuras en el sitio indicado H.


% Paquetes de imágenes
\usepackage{graphicx}       % Paquete para añadir imágenes
\usepackage{transparent}    % Para manejar la opacidad de las figuras

% Paquete para usar colores
\usepackage[dvipsnames]{xcolor}
\usepackage{pagecolor}      % Para cambiar el color de la página

% Habilita tamaños de fuente mayores
\usepackage{fix-cm}


%------------------------------------------------------------------------

% Paquetes de matemáticas
\usepackage{mathtools, amsfonts, amssymb, mathrsfs}
\usepackage[makeroom]{cancel}     % Simplificar tachando
\usepackage{polynom}    % Divisiones y Ruffini
\usepackage{units} % Para poner fracciones diagonales con \nicefrac

\usepackage{pgfplots}   %Representar funciones
\pgfplotsset{compat=1.18}  % Versión 1.18

\usepackage{tikz-cd}    % Para usar diagramas de composiciones
\usetikzlibrary{calc}   % Para usar cálculo de coordenadas en tikz

%Definición de teoremas, etc.
\usepackage{amsthm}
%\swapnumbers   % Intercambia la posición del texto y de la numeración

\theoremstyle{plain}

\makeatletter
\@ifclassloaded{article}{
  \newtheorem{teo}{Teorema}[section]
}{
  \newtheorem{teo}{Teorema}[chapter]  % Se resetea en cada chapter
}
\makeatother

\newtheorem{coro}{Corolario}[teo]           % Se resetea en cada teorema
\newtheorem{prop}[teo]{Proposición}         % Usa el mismo contador que teorema
\newtheorem{lema}[teo]{Lema}                % Usa el mismo contador que teorema

\theoremstyle{remark}
\newtheorem*{observacion}{Observación}

\theoremstyle{definition}

\makeatletter
\@ifclassloaded{article}{
  \newtheorem{definicion}{Definición} [section]     % Se resetea en cada chapter
}{
  \newtheorem{definicion}{Definición} [chapter]     % Se resetea en cada chapter
}
\makeatother

\newtheorem*{notacion}{Notación}
\newtheorem*{ejemplo}{Ejemplo}
\newtheorem*{ejercicio*}{Ejercicio}             % No numerado
\newtheorem{ejercicio}{Ejercicio} [section]     % Se resetea en cada section


% Modificar el formato de la numeración del teorema "ejercicio"
\renewcommand{\theejercicio}{%
  \ifnum\value{section}=0 % Si no se ha iniciado ninguna sección
    \arabic{ejercicio}% Solo mostrar el número de ejercicio
  \else
    \thesection.\arabic{ejercicio}% Mostrar número de sección y número de ejercicio
  \fi
}


% \renewcommand\qedsymbol{$\blacksquare$}         % Cambiar símbolo QED
%------------------------------------------------------------------------

% Paquetes para encabezados
\usepackage{fancyhdr}
\pagestyle{fancy}
\fancyhf{}

\newcommand{\helv}{ % Modificación tamaño de letra
\fontfamily{}\fontsize{12}{12}\selectfont}
\setlength{\headheight}{15pt} % Amplía el tamaño del índice


%\usepackage{lastpage}   % Referenciar última pag   \pageref{LastPage}
\fancyfoot[C]{\thepage}

%------------------------------------------------------------------------

% Conseguir que no ponga "Capítulo 1". Sino solo "1."
\makeatletter
\@ifclassloaded{book}{
  \renewcommand{\chaptermark}[1]{\markboth{\thechapter.\ #1}{}} % En el encabezado
    
  \renewcommand{\@makechapterhead}[1]{%
  \vspace*{50\p@}%
  {\parindent \z@ \raggedright \normalfont
    \ifnum \c@secnumdepth >\m@ne
      \huge\bfseries \thechapter.\hspace{1em}\ignorespaces
    \fi
    \interlinepenalty\@M
    \Huge \bfseries #1\par\nobreak
    \vskip 40\p@
  }}
}
\makeatother

%------------------------------------------------------------------------
% Paquetes de cógido
\usepackage{minted}
\renewcommand\listingscaption{Código fuente}

\usepackage{fancyvrb}
% Personaliza el tamaño de los números de línea
\renewcommand{\theFancyVerbLine}{\small\arabic{FancyVerbLine}}

% Estilo para C++
\newminted{cpp}{
    frame=lines,
    framesep=2mm,
    baselinestretch=1.2,
    linenos,
    escapeinside=||
}



\usepackage{listings} % Para incluir código desde un archivo

\renewcommand\lstlistingname{Código Fuente}
\renewcommand\lstlistlistingname{Índice de Códigos Fuente}

% Definir colores
\definecolor{vscodepurple}{rgb}{0.5,0,0.5}
\definecolor{vscodeblue}{rgb}{0,0,0.8}
\definecolor{vscodegreen}{rgb}{0,0.5,0}
\definecolor{vscodegray}{rgb}{0.5,0.5,0.5}
\definecolor{vscodebackground}{rgb}{0.97,0.97,0.97}
\definecolor{vscodelightgray}{rgb}{0.9,0.9,0.9}

% Configuración para el estilo de C similar a VSCode
\lstdefinestyle{vscode_C}{
  backgroundcolor=\color{vscodebackground},
  commentstyle=\color{vscodegreen},
  keywordstyle=\color{vscodeblue},
  numberstyle=\tiny\color{vscodegray},
  stringstyle=\color{vscodepurple},
  basicstyle=\scriptsize\ttfamily,
  breakatwhitespace=false,
  breaklines=true,
  captionpos=b,
  keepspaces=true,
  numbers=left,
  numbersep=5pt,
  showspaces=false,
  showstringspaces=false,
  showtabs=false,
  tabsize=2,
  frame=tb,
  framerule=0pt,
  aboveskip=10pt,
  belowskip=10pt,
  xleftmargin=10pt,
  xrightmargin=10pt,
  framexleftmargin=10pt,
  framexrightmargin=10pt,
  framesep=0pt,
  rulecolor=\color{vscodelightgray},
  backgroundcolor=\color{vscodebackground},
}

%------------------------------------------------------------------------

% Comandos definidos
\newcommand{\bb}[1]{\mathbb{#1}}
\newcommand{\cc}[1]{\mathcal{#1}}

% I prefer the slanted \leq
\let\oldleq\leq % save them in case they're every wanted
\let\oldgeq\geq
\renewcommand{\leq}{\leqslant}
\renewcommand{\geq}{\geqslant}

% Si y solo si
\newcommand{\sii}{\iff}

% Letras griegas
\newcommand{\eps}{\epsilon}
\newcommand{\veps}{\varepsilon}
\newcommand{\lm}{\lambda}

\newcommand{\ol}{\overline}
\newcommand{\ul}{\underline}
\newcommand{\wt}{\widetilde}
\newcommand{\wh}{\widehat}

\let\oldvec\vec
\renewcommand{\vec}{\overrightarrow}

% Derivadas parciales
\newcommand{\del}[2]{\frac{\partial #1}{\partial #2}}
\newcommand{\Del}[3]{\frac{\partial^{#1} #2}{\partial^{#1} #3}}
\newcommand{\deld}[2]{\dfrac{\partial #1}{\partial #2}}
\newcommand{\Deld}[3]{\dfrac{\partial^{#1} #2}{\partial^{#1} #3}}


\newcommand{\AstIg}{\stackrel{(\ast)}{=}}
\newcommand{\Hop}{\stackrel{L'H\hat{o}pital}{=}}

\newcommand{\red}[1]{{\color{red}#1}} % Para integrales, destacar los cambios.

% Método de integración
\newcommand{\MetInt}[2]{
    \left[\begin{array}{c}
        #1 \\ #2
    \end{array}\right]
}

% Declarar aplicaciones
% 1. Nombre aplicación
% 2. Dominio
% 3. Codominio
% 4. Variable
% 5. Imagen de la variable
\newcommand{\Func}[5]{
    \begin{equation*}
        \begin{array}{rrll}
            #1:& #2 & \longrightarrow & #3\\
               & #4 & \longmapsto & #5
        \end{array}
    \end{equation*}
}

%------------------------------------------------------------------------

% Para poder incluir árboles
\usepackage{forest}
\usepackage{booktabs}

% Para poder añadir autómatas
% https://www3.nd.edu/~kogge/courses/cse30151-fa17/Public/other/tikz_tutorial.pdf
\usetikzlibrary{automata} %, positioning, arrows}
\tikzset{
    -Stealth,
    node distance=3cm, % specifies the minimum distance between two nodes. Change if necessary.
    every state/.style={thick, fill=gray!10, shape=ellipse}, % sets the properties for each ’state’ node
    initial text=$ $, % sets the text that appears on the start arrow
    % Un tipo de nodo, que es error, que lo pone rojo
    error/.style={thick, fill=red!20},
}

\begin{document}

    % 1. Foto de fondo
    % 2. Título
    % 3. Encabezado Izquierdo
    % 4. Color de fondo
    % 5. Coord x del titulo
    % 6. Coord y del titulo
    % 7. Fecha

    
    % 1. Foto de fondo
% 2. Título
% 3. Encabezado Izquierdo
% 4. Color de fondo
% 5. Coord x del titulo
% 6. Coord y del titulo
% 7. Fecha

\newcommand{\portada}[7]{

    \portadaBase{#1}{#2}{#3}{#4}{#5}{#6}{#7}
    \portadaBook{#1}{#2}{#3}{#4}{#5}{#6}{#7}
}

\newcommand{\portadaExamen}[7]{

    \portadaBase{#1}{#2}{#3}{#4}{#5}{#6}{#7}
    \portadaArticle{#1}{#2}{#3}{#4}{#5}{#6}{#7}
}




\newcommand{\portadaBase}[7]{

    % Tiene la portada principal y la licencia Creative Commons
    
    % 1. Foto de fondo
    % 2. Título
    % 3. Encabezado Izquierdo
    % 4. Color de fondo
    % 5. Coord x del titulo
    % 6. Coord y del titulo
    % 7. Fecha
    
    
    \thispagestyle{empty}               % Sin encabezado ni pie de página
    \newgeometry{margin=0cm}        % Márgenes nulos para la primera página
    
    
    % Encabezado
    \fancyhead[L]{\helv #3}
    \fancyhead[R]{\helv \nouppercase{\leftmark}}
    
    
    \pagecolor{#4}        % Color de fondo para la portada
    
    \begin{figure}[p]
        \centering
        \transparent{0.3}           % Opacidad del 30% para la imagen
        
        \includegraphics[width=\paperwidth, keepaspectratio]{assets/#1}
    
        \begin{tikzpicture}[remember picture, overlay]
            \node[anchor=north west, text=white, opacity=1, font=\fontsize{60}{90}\selectfont\bfseries\sffamily, align=left] at (#5, #6) {#2};
            
            \node[anchor=south east, text=white, opacity=1, font=\fontsize{12}{18}\selectfont\sffamily, align=right] at (9.7, 3) {\textbf{\href{https://losdeldgiim.github.io/}{Los Del DGIIM}}};
            
            \node[anchor=south east, text=white, opacity=1, font=\fontsize{12}{15}\selectfont\sffamily, align=right] at (9.7, 1.8) {Doble Grado en Ingeniería Informática y Matemáticas\\Universidad de Granada};
        \end{tikzpicture}
    \end{figure}
    
    
    \restoregeometry        % Restaurar márgenes normales para las páginas subsiguientes
    \pagecolor{white}       % Restaurar el color de página
    
    
    \newpage
    \thispagestyle{empty}               % Sin encabezado ni pie de página
    \begin{tikzpicture}[remember picture, overlay]
        \node[anchor=south west, inner sep=3cm] at (current page.south west) {
            \begin{minipage}{0.5\paperwidth}
                \href{https://creativecommons.org/licenses/by-nc-nd/4.0/}{
                    \includegraphics[height=2cm]{assets/Licencia.png}
                }\vspace{1cm}\\
                Esta obra está bajo una
                \href{https://creativecommons.org/licenses/by-nc-nd/4.0/}{
                    Licencia Creative Commons Atribución-NoComercial-SinDerivadas 4.0 Internacional (CC BY-NC-ND 4.0).
                }\\
    
                Eres libre de compartir y redistribuir el contenido de esta obra en cualquier medio o formato, siempre y cuando des el crédito adecuado a los autores originales y no persigas fines comerciales. 
            \end{minipage}
        };
    \end{tikzpicture}
    
    
    
    % 1. Foto de fondo
    % 2. Título
    % 3. Encabezado Izquierdo
    % 4. Color de fondo
    % 5. Coord x del titulo
    % 6. Coord y del titulo
    % 7. Fecha


}


\newcommand{\portadaBook}[7]{

    % 1. Foto de fondo
    % 2. Título
    % 3. Encabezado Izquierdo
    % 4. Color de fondo
    % 5. Coord x del titulo
    % 6. Coord y del titulo
    % 7. Fecha

    % Personaliza el formato del título
    \pretitle{\begin{center}\bfseries\fontsize{42}{56}\selectfont}
    \posttitle{\par\end{center}\vspace{2em}}
    
    % Personaliza el formato del autor
    \preauthor{\begin{center}\Large}
    \postauthor{\par\end{center}\vfill}
    
    % Personaliza el formato de la fecha
    \predate{\begin{center}\huge}
    \postdate{\par\end{center}\vspace{2em}}
    
    \title{#2}
    \author{\href{https://losdeldgiim.github.io/}{Los Del DGIIM}}
    \date{Granada, #7}
    \maketitle
    
    \tableofcontents
}




\newcommand{\portadaArticle}[7]{

    % 1. Foto de fondo
    % 2. Título
    % 3. Encabezado Izquierdo
    % 4. Color de fondo
    % 5. Coord x del titulo
    % 6. Coord y del titulo
    % 7. Fecha

    % Personaliza el formato del título
    \pretitle{\begin{center}\bfseries\fontsize{42}{56}\selectfont}
    \posttitle{\par\end{center}\vspace{2em}}
    
    % Personaliza el formato del autor
    \preauthor{\begin{center}\Large}
    \postauthor{\par\end{center}\vspace{3em}}
    
    % Personaliza el formato de la fecha
    \predate{\begin{center}\huge}
    \postdate{\par\end{center}\vspace{5em}}
    
    \title{#2}
    \author{\href{https://losdeldgiim.github.io/}{Los Del DGIIM}}
    \date{Granada, #7}
    \thispagestyle{empty}               % Sin encabezado ni pie de página
    \maketitle
    \vfill
}
    \portadaExamen{etsiitA4.jpg}{Modelos de\\Computación\\Examen IV}{MC. Examen IV}{MidnightBlue}{-8}{28}{2024-2025}{Arturo Olivares Martos}

    \begin{description}
        \item[Asignatura] Modelos de Computación
        \item[Curso Académico] 2022-23.
        \item[Grado] Doble Grado en Ingeniería Informática y Matemáticas.
        %\item[Grupo] B.
        %\item[Profesor] Rafael Ortega Ríos.
        \item[Descripción] Parcial Tema 2.
        %\item[Fecha] 22 de marzo de 2018.
        %\item[Duración] 60 minutos.    
    \end{description}
    \newpage
    
    \begin{ejercicio} \label{ej:1}
        Encuentra un AFD que acepte el lenguaje descrito por la expresión regular:
        \begin{equation*}
            a+ac(a+b)^* + c(a+b+c)^*.
        \end{equation*}

        El AFND con transiciones nulas que optenemos usando el algoritmo viene descrito en la Figura~\ref{fig:ej1_AFND}.
        \begin{figure}[h]
            \centering
            \begin{tikzpicture}
                \node[state, initial] (q0) {$q_0$};
                \node[state, above right of=q0] (q01) {$q_{01}$};
                \node[state, right of=q0, xshift=-2em] (q02) {$q_{02}$};
                \node[state, below right of=q0] (q03) {$q_{03}$};
                \node[state, right of=q01, accepting] (q11) {$q_{11}$};
                \node[state, right of=q02] (q12) {$q_{12}$};
                \node[state, right of=q12, accepting] (q22) {$q_{22}$};
                \node[state, right of=q22] (q32) {$q_{32}$};
                \node[state, above right of=q32, accepting] (q032) {$q_{032}$};
                \node[state, below right of=q32, accepting] (q132) {$q_{132}$};
                \node[state, right of=q03, accepting] (q13) {$q_{13}$};
                \node[state, below of=q13] (q23) {$q_{23}$};
                \node[state, below left of=q23, accepting] (q023) {$q_{023}$};
                \node[state, below of=q23, accepting, yshift=1em] (q123) {$q_{123}$};
                \node[state, below right of=q23, accepting] (q223) {$q_{223}$};

                \draw   (q0) edge[above] node{$\varepsilon$} (q01)
                        (q0) edge[above] node{$\varepsilon$} (q02)
                        (q0) edge[above] node{$\varepsilon$} (q03)
                        (q01) edge[above] node{$a$} (q11)
                        (q02) edge[above] node{$a$} (q12)
                        (q12) edge[above] node{$c$} (q22)
                        (q22) edge[above] node{$\veps$} (q32)
                        (q32) edge[above] node{$a$} (q032)
                        (q32) edge[above] node{$b$} (q132)
                        (q032) edge[above, bend right=60] node{$\veps$} (q32)
                        (q132) edge[below, bend left=60] node{$\veps$} (q32)
                        (q03) edge[above] node{$c$} (q13)
                        (q13) edge[right] node{$\veps$} (q23)
                        (q23) edge[above] node{$a$} (q023)
                        (q23) edge[right] node{$b$} (q123)
                        (q23) edge[above] node{$c$} (q223)
                        (q023) edge[above, bend left=60] node{$\veps$} (q23)
                        (q123) edge[left, bend left] node[pos=0.1]{$\veps$} (q23)
                        (q223) edge[above, bend right=60] node{$\veps$} (q23);
            \end{tikzpicture}
            \caption{AFND con transiciones nulas que acepta el lenguaje del Ejercicio~\ref{ej:1}.}
            \label{fig:ej1_AFND}
        \end{figure}

        No obstante, este AFND es demasiado complejo y tiene demasiados estados, por lo que lo minimizamos antes de convertirlo en un AFD. El AFND minimizado se muestra en la Figura~\ref{fig:ej1_AFND_min}.
        \begin{figure}[h]
            \centering
            \begin{tikzpicture}
                \node[state, initial] (q0) {$q_0$};
                \node[state, right of=q0, xshift=-2em, accepting] (q12) {$q_{12}$};
                \node[state, right of=q12, accepting] (q22) {$q_{22}$};
                \node[state, below right of=q0, accepting] (q13) {$q_{13}$};

                \draw   (q0) edge[above] node{$a$} (q12)
                        (q12) edge[above] node{$c$} (q22)
                        (q22) edge[loop above] node{$a,b$} (q22)
                        (q0) edge[above] node{$c$} (q13)
                        (q13) edge[loop right] node{$a,b,c$} (q13);
            \end{tikzpicture}
            \caption{AFND minimizado que acepta el lenguaje del Ejercicio~\ref{ej:1}.}
            \label{fig:ej1_AFND_min}
        \end{figure}

        Finalmente, convertimos el AFND minimizado en un AFD, que se muestra en la Figura~\ref{fig:ej1_AFD}.
        \begin{figure}[h]
            \centering
            \begin{tikzpicture}
                \node[state, initial] (q0) {$q_0$};
                \node[state, right of=q0, xshift=-2.1em, accepting] (q12) {$q_{12}$};
                \node[state, right of=q12, accepting] (q22) {$q_{22}$};
                \node[state, below right of=q0, accepting] (q13) {$q_{13}$};
                \node[state, above right of=q0, error] (E) {$E$};

                \draw   (q0) edge[above] node{$a$} (q12)
                        (q12) edge[above] node{$c$} (q22)
                        (q22) edge[loop above] node{$a,b$} (q22)
                        (q0) edge[above] node{$c$} (q13)
                        (q13) edge[loop right] node{$a,b,c$} (q13);

                % Conexiones para el estado de error
                \draw[opacity=0.6]   (q0) edge[above] node{$b$} (E)
                        (q12) edge[right] node{$a,b$} (E)
                        (q22) edge[above] node{$c$} (E)
                        (E) edge[loop right] node{$a,b,c$} (E);
            \end{tikzpicture}
            \caption{AFD que acepta el lenguaje del Ejercicio~\ref{ej:1}.}
            \label{fig:ej1_AFD}
        \end{figure}

        

        
    \end{ejercicio}

    \begin{ejercicio} \label{ej:2}
        Considera el lenguaje de todas las palabras en las que toda subcadena de $1$'s de longitud mayor o igual a $2$ está precedida de una subcadena de $0$'s de cualquier longitud mayor o igual a $3$. Encuentra un autómata finito, de cualquier tipo, que acepte este lenguaje.\\

        Consideramos los siguientes estados:
        \begin{itemize}
            \item $q_0$: Estado inicial.
            \item $q_1$: Hemos leído un $1$ que no está precedido por $3$ o más $0$'s.
            \item $q_2$: Hemos leído solo un $0$ consecutivo.
            \item $q_3$: Hemos leído dos $0$'s consecutivos.
            \item $q_4$: Hemos leído tres $0$'s consecutivos.
            \item $q_5$: He leído un $1$ que está precedido por $3$ o más $0$'s, por lo que puedo introducir más $1$'s.
        \end{itemize}

        El AFD que acepta el lenguaje descrito se muestra en la Figura~\ref{fig:ej2_AFD}.
        \begin{figure}[h]
            \centering
            \begin{tikzpicture}
                \node[state, initial, accepting] (q0) {$q_0$};
                \node[state, right of=q0, accepting] (q1) {$q_1$};
                \node[state, right of=q1, accepting] (q2) {$q_2$};
                \node[state, right of=q2, accepting] (q3) {$q_3$};
                \node[state, right of=q3, accepting] (q4) {$q_4$};
                \node[state, right of=q4, accepting] (q5) {$q_5$};
                \node[state, above of=q1, error] (E) {$E$};

                \draw   (q0) edge[above] node{$1$} (q1)
                        (q0) edge[bend right, below] node{$0$} (q2)
                        (q1) edge[below] node{$0$} (q2)
                        (q1) edge[left] node{$1$} (E)
                        (q2) edge[below] node{$0$} (q3)
                        (q2) edge[bend right, above] node[pos=0.1]{$1$} (q1)
                        (q3) edge[above] node{$0$} (q4)
                        (q3) edge[bend right, above] node{$1$} (q1)
                        (q4) edge[loop above] node{$0$} (q4)
                        (q4) edge[above] node{$1$} (q5)
                        (q5) edge[loop above] node{$1$} (q5)
                        (q5) edge[bend left, above] node{$0$} (q2)
                        (E) edge[loop right] node{$0,1$} (E);
            \end{tikzpicture}
            \caption{AFD que acepta el lenguaje del Ejercicio~\ref{ej:2}.}
            \label{fig:ej2_AFD}
        \end{figure}
    \end{ejercicio}

    \begin{ejercicio} \label{ej:3}
        Considera el lenguaje $L\subset \{a,b\}^*$ formada por las palabras en las que las $a$'s y las $b$'s están siempre alternadas. Ejemplos de palabras del lenguaje son:
        \begin{equation*}
            \veps,a,b,ab,ba,aba,bab,abab,baba \in L.
        \end{equation*}
        Ejemplos de palabras que no están en el lenguaje son:
        \begin{equation*}
            bb,aa,abba \notin L.
        \end{equation*}
        Encontrar una expresión regular que genere el lenguaje $L$.\\

        Obtenemos en primer lugar un AFD que acepte el lenguaje $L$. Este tiene los siguientes estados:
        \begin{itemize}
            \item $q_0$: Estado inicial.
            \item $q_A$: Estado en el que la última letra leída es una $a$.
            \item $q_B$: Estado en el que la última letra leída es una $b$.
        \end{itemize}

        Este AFD se muestra en la Figura~\ref{fig:ej3_AFD}.
        \begin{figure}[h]
            \centering
            \begin{tikzpicture}
                \node[state, initial, accepting] (q0) {$q_0$};
                \node[state, right of=q0, accepting, yshift=3em] (qA) {$q_A$};
                \node[state, right of=q0, accepting, yshift=-3em] (qB) {$q_B$};
                \node[state, right of=qB, error, yshift=3em] (E) {$E$};

                \draw   (q0) edge[above] node{$a$} (qA)
                        (q0) edge[below] node{$b$} (qB)
                        (qA) edge[left, bend right] node{$b$} (qB)
                        (qB) edge[right, bend right] node{$a$} (qA)
                        (qA) edge[above] node{$a$} (E)
                        (qB) edge[below] node{$b$} (E)
                        (E) edge[loop right] node{$a,b$} (E);
            \end{tikzpicture}
            \caption{AFD que acepta el lenguaje del Ejercicio~\ref{ej:3}.}
            \label{fig:ej3_AFD}
        \end{figure}

        Para obtener la expresión regular que genera el lenguaje $L$, planteamos el sistema de ecuaciones que se muestra a continuación:
        \begin{align*}
            q_0 &= aq_A + bq_B + \veps,\\
            q_A &= aE + bq_B + \veps,\\
            q_B &= aq_A + bE + \veps,\\
            E &= aE + bE = (a+b)E + \emptyset
        \end{align*}
        Por el lema de Arden, tenemos que $E=(a+b)^*\emptyset = \emptyset$. Sustituyendo en las ecuaciones anteriores, obtenemos:
        \begin{align*}
            q_0 &= aq_A + bq_B + \veps,\\
            q_A &= bq_B + \veps,\\
            q_B &= aq_A + \veps.
        \end{align*}

        Sustituyendo en valor de $q_B$, obtenemos:
        \begin{align*}
            q_0 &= aq_A + b(aq_A + \veps) + \veps,\\
            q_A &= b(aq_A + \veps) + \veps = baq_A + b + \veps
        \end{align*}

        Usando de nuevo el lema de Arden, tenemos que $q_A = (ba)^*(b+\veps)$. Sustituyendo en la ecuación de $q_0$, obtenemos:
        \begin{align*}
            q_0 &= a(ba)^*(b+\veps) + b(a(ba)^*(b+\veps) + \veps) + \veps
            =\\&= a(ba)^*(b+\veps) + (ba)^+(b+\veps) + b + \veps
        \end{align*}

        Por lo tanto, la expresión regular que genera el lenguaje $L$ es:
        \begin{equation*}
            a(ba)^*(b+\veps) + (ba)^+(b+\veps) + b + \veps.
        \end{equation*}
    \end{ejercicio}
\end{document}
